\section{Conclusion}

This master's thesis has developed a continuous bilevel optimization framework for airline hub-and-spoke network design that simultaneously addresses three modeling challenges: the combinatorial nature of facility activation, the nonconvexity of logit demand capture, and the calibration of market-specific priorities against the airline's true profit objective.

\paragraph{Summary of contributions.}
The first contribution is a continuous lower-level model in which reweighted $\ell_1$ penalties replace the binary activation variables for airports, hubs, and service capacities, while entropy regularization replaces the explicit logit demand-capture equation. The resulting lower-level problem is convex for fixed market weights, amenable to off-the-shelf interior-point solvers, and provably recovers logit market shares as its optimality conditions through a convex-conjugate construction \cite{Oppenheim1993,RealRojas2026TRB}. Sparse deployment is promoted through a majorization-minimization outer loop that iteratively tightens the convex surrogate of the $\ell_0$ fixed-cost structure around the current design.

The second contribution is the bilevel layer on top of this continuous model. The upper level selects market-specific weight coefficients $(\widehat{\alpha},\widehat{\beta})$ to maximize the airline's true profit, while the lower level returns the network design induced by those weights. This hierarchical structure allows the framework to calibrate market priorities jointly with the network design, instead of relying on exogenously chosen parameters. The upper-level objective is the real profit function, independent of the lower-level weighting, so the bilevel formulation provides a principled way of steering a continuous convex approximation toward economically meaningful solutions.

The third contribution is the penalty-based solution algorithm. Following the value-function penalty approach of \cite{Ye1997Exact,Jiang2024Primal}, the bilevel optimality condition is relaxed into a soft penalty added to the upper-level objective, producing a single-level surrogate that is tractable by gradient-based methods. The inner penalty loop alternates lower-level and penalized solves to update the market coefficients via projected steps, while the outer majorization-minimization loop progressively sharpens the sparse deployment approximation.

\paragraph{Computational findings.}
The methodology has been validated on Spanish domestic instances of four, six, and eight airports, and on a realistic Madrid-centred international network of twenty-five airports with Iberia-style coverage. On the small Spanish instances, the bilevel scheme solves in the second-to-minute range while the mixed-integer baseline approaches or exceeds a two-hour time limit, representing speedups of one to three orders of magnitude. The profit gap relative to the mixed-integer optimum is small for several budget and step-size configurations, confirming that the continuous approximation is economically meaningful and not merely fast. The hub structure recovered by the bilevel scheme is qualitatively consistent with the mixed-integer solution in all tested instances, with MAD acting as the dominant hub and secondary hubs emerging only when the budget is sufficient to support them. On the international instance, the mixed-integer baseline fails to produce a feasible incumbent within the time limit, whereas the bilevel scheme consistently returns sparse network designs with a concentrated set of European hubs and a coherent allocation of transatlantic services.

\paragraph{Limitations.}
Several limitations of the present work should be acknowledged. First, the penalty-based bilevel algorithm is not guaranteed to converge to a globally optimal bilevel solution; the theoretical guarantees cover only stationarity neighborhoods, and the practical quality of the solution depends on the choice of step sizes $\mu_\alpha,\mu_\beta$ and penalty weight $\gamma$. Second, the continuous lower-level model approximates fixed activation costs through a smooth surrogate that is only exact in the limit of the majorization-minimization iterations; for finite iteration counts, some residual activity may remain at nominally inactive facilities. Third, the current demand model treats fares as exogenous parameters and does not consider endogenous pricing decisions by the entrant or strategic responses by incumbent carriers, which limits the applicability of the framework to settings where competitive equilibria are relevant.

\paragraph{Future work.}
Several directions of future research follow naturally from this work. On the algorithmic side, the convergence guarantees of the penalty loop could be strengthened by adopting double-loop penalty schemes with adaptive penalty sequences \cite{LuMei2024}, or by replacing the gradient-like market-coefficient update with a proper implicit-gradient step obtained through the implicit function theorem. On the modeling side, the framework could be extended to incorporate endogenous pricing by the entrant as an additional upper-level decision, turning the model into a full Stackelberg pricing-and-design bilevel program. A further natural extension is to include multiple competing airlines as separate lower-level players, resulting in a multi-follower bilevel model closer to a market equilibrium formulation. Finally, the reweighted $\ell_1$ approximation of fixed costs could be replaced by tighter DC (difference-of-convex) surrogates or by mixed-integer rounding strategies applied to the continuous solution, potentially reducing the gap to the exact mixed-integer optimum at moderate additional computational cost.

In summary, this master's thesis demonstrates that the combination of entropy-based demand modeling, sparsity-inducing continuous activation, and value-function-penalty bilevel calibration provides a practical and computationally viable route to airline hub-and-spoke network design at scales that are beyond the reach of exact mixed-integer bilevel formulations, while maintaining an explicit and interpretable link between passenger behavior, network topology, and airline profitability.
